% Ч1.2 — Среднее арифметическое
% Старая методичка: стр. 216 (старое задание 19)
% Источники: «ВСЯ ТЕОРИЯ ЕГЭ» (стр. 216); «Задача 20: необходимая теория», § 14;
% методичка по заданию 11 (средняя скорость). Стиль — стиль-изложения.md

\begin{tema}{Ч1.2}
\opr{Средним арифметическим} чисел $a_1, a_2, \dots, a_n$ называется число
\[
  m=\frac{a_1+a_2+\dots+a_n}{n}.
\]
Отсюда сумма чисел равна среднему арифметическому, умноженному на их количество:
\[
  a_1+a_2+\dots+a_n=m\cdot n.
\]

\begin{primer}
Среднее арифметическое четырёх оценок равно 4. Какой должна быть пятая оценка,
чтобы среднее стало равно 4,2?
Решение: сумма четырёх оценок $4\cdot4=16$, а пяти должна быть $4{,}2\cdot5=21$.
Значит, пятая оценка $21-16=5$. Ответ: 5.
\end{primer}

\opr{Примечание.} Математическое ожидание (\tref{В3.2}) — это тоже среднее,
только «с весами»: каждое значение умножается на свою вероятность,
и результаты складываются.

\begin{oshibka}
Средняя скорость в общем случае — \emph{не} среднее арифметическое скоростей,
а весь путь, делённый на всё время: $v_\text{ср}=\frac{S_\text{весь}}{t_\text{всё}}$ (\tref{М1.7}).
Совпадают они, только если на каждой скорости ехали одинаковое время.
Половину пути со скоростью 60 км/ч и половину — со скоростью 40 км/ч:
средняя скорость 48 км/ч, а не 50.
\end{oshibka}
\end{tema}
